# EMRA Robot Pilot Evaluation Instruments

This repository presents the complete evaluation instruments used in the autism-center pilot study of the **Educational Multimodal Robot Assistant (EMRA)**.

The evaluation package consists of:

- **F1:** Child Session Observation Form
- **F2:** Child Preference Form
- **F3:** Educator Post-Pilot Questionnaire
- **Automatically recorded EMRA system fields**

---

# F1: Child Session Observation Form

This form is completed by the observer during and immediately after the session. One form is completed for each child session.

## Session Information

| Field | Information |
|---|---|
| Child ID | |
| Session Number | |
| Language | ☐ Arabic  ☐ English  ☐ Auto |
| Start Time | |
| End Time | |

## Part 1: Educational Activities

### Prompt-Level Scale

| Level | Definition |
|:---:|---|
| 0 | Independent. The child completed the activity with no help beyond the robot’s first instruction. |
| 1 | The robot repeated or rephrased the instruction, and the child then responded. |
| 2 | The robot gave an additional cue, such as an on-screen highlight, or the child pressed **Help**. |
| 3 | An adult gave a verbal prompt. |
| 4 | An adult gave physical guidance, only where this is already part of normal center practice. |
| N/A | The activity was opened but not attempted. |

### Educational Activity Record

| Activity | Prompt Level | Completed | Loss of Attention | Play Again | Notes |
|---|:---:|:---:|:---:|:---:|---|
| Body Parts | | | | | |
| Puzzle | | | | | |
| LEGO | | | | | |
| Painting | | | | | |
| Colors | | | | | |
| Numbers | | | | | |
| Writing | | | | | |
| Reading | | | | | |
| Letters | | | | | |

**Recording options:**

- **Prompt Level:** 0–4
- **Completed:** Yes / Partial / No
- **Loss of Attention:** Never / Sometimes / Often / Very often

## Part 2: Talk Session

Complete this part only when the child uses the **Talk** session. Skip it when the child uses only the educational activities.

### 1. How often did the child give a clear response after EMRA spoke?

- ○ Never
- ○ Rarely
- ○ Sometimes
- ○ Often

**Notes:**  
<br><br>

### 2. How often did the child speak to EMRA without being asked?

- ○ Never
- ○ Rarely
- ○ Sometimes
- ○ Often

**Notes:**  
<br><br>

### 3. How often did the child look at EMRA’s face screen?

- ○ Never
- ○ Rarely
- ○ Sometimes
- ○ Often

**Notes:**  
<br><br>

### 4. When a physical cue occurred while the child’s attention was away, did the child redirect attention toward EMRA or the activity?

- ○ Yes
- ○ No
- ○ N/A

**Notes:**  
<br><br>

### 5. Which gesture affected the child more?

- ☐ Nod
- ☐ Shake
- ☐ Handshake
- ☐ Wave
- ☐ Hug
- ☐ Clap

**Notes:**  
<br><br>

### 6. Did the child show the robot an object or drawing?

- ○ Yes
- ○ No

**Notes:**  
<br><br>

### 7. How often did an adult have to explain or repeat what EMRA said because the child did not understand or respond directly?

- ○ Never
- ○ Rarely
- ○ Sometimes
- ○ Often

**Notes:**  
<br><br>

### 8. How often did EMRA reply in the wrong language?

- ○ Never
- ○ Rarely
- ○ Sometimes
- ○ Often

**Notes:**  
<br><br>

## Part 3: Gesture-by-Gesture Response

For every gesture performed by EMRA, record the child’s observable response. Leave gestures that did not occur blank.

The system automatically records which gesture was performed and why it was selected.

### Child-Response Codes

| Code | Child’s Observable Response |
|:---:|---|
| R0 | No observable response. |
| R1 | Looked toward EMRA. |
| R2 | Looked toward the moving arm, hand, or head. |
| R3 | Responded verbally. |
| R4 | Imitated the gesture. |
| R5 | Moved closer to EMRA or touched it. |
| R6 | Moved away or showed discomfort. |

### Gesture-Response Record

| Gesture | Frequency | Response Codes | Notes |
|---|:---:|:---:|---|
| Nod | | | |
| Shake | | | |
| Handshake | | | |
| Wave | | | |
| Hug | | | |
| Clap | | | |

> A single gesture may produce more than one response code during a session. Record all observed response codes.

## Part 4: Global Session Ratings

Select one score from **1 to 5** for each measure.

| Measure | 1 | 2 | 3 | 4 | 5 |
|---|---|---|---|---|---|
| **Engagement** | No participation | Adult-directed only | Some participation | Participated most of the session | Highly engaged |
| **Attention** | Almost none | Brief | About half | Most of the session | Sustained throughout |
| **Affect** | Distressed | Uncomfortable | Neutral | Mostly positive | Clearly positive |
| **Cooperation** | Did not follow | Needed adult repetition | Followed some | Followed most | Followed consistently |
| **Willingness to continue** | Wanted to stop | Wanted to stop early | Neutral | Willing to continue | Clearly wanted to continue |
| **Emotional appropriateness** | Clearly mismatched | Often mismatched | Mixed or unclear | Mostly appropriate | Clearly appropriate |

> Emotional appropriateness refers to the spoken reply and face-screen expression considered together.

## Part 5: Event Checklist

### 1. Child reacted negatively to the robot’s sound

- ○ Never
- ○ Rarely
- ○ Sometimes
- ○ Often

**Notes:**  
<br><br>

### 2. Child reacted negatively to the robot’s movement

- ○ Never
- ○ Rarely
- ○ Sometimes
- ○ Often

**Notes:**  
<br><br>

### 3. Child showed distress

- ○ Never
- ○ Rarely
- ○ Sometimes
- ○ Often

**Notes:**  
<br><br>

### 4. Child moved away from or refused EMRA

- ○ Never
- ○ Rarely
- ○ Sometimes
- ○ Often

**Notes:**  
<br><br>

### 5. EMRA did not respond or gave an unrelated reply

- ○ Never
- ○ Rarely
- ○ Sometimes
- ○ Often

**Notes:**  
<br><br>

### 6. EMRA’s facial expression did not match the child’s observed state

- ○ Never
- ○ Rarely
- ○ Sometimes
- ○ Often

**Notes:**  
<br><br>

### 7. Noticeable response delay or technical failure

- ○ Never
- ○ Rarely
- ○ Sometimes
- ○ Often

**Notes:**  
<br><br>

### 8. Session stopped early

- ○ No
- ○ Child-related
- ○ Technical

**Notes:**  
<br><br>

### 9. Other important event

**Description:**  
<br><br><br>

---

# F2: Child Preference Form

Complete this form immediately after each session with educator assistance.

## Session Information

| Field | Information |
|---|---|
| Child ID | |
| Session Number | |
| Date | |

## Instructions

**Ask the child to point to the face that shows what they think.**

## Q1. How was playing with EMRA today?

| 😀 Happy | 😐 Neutral | ☹️ Sad |
|:---:|:---:|:---:|
| ○ | ○ | ○ |

## Q2. How was talking with EMRA today?

Mark **N/A** if the Talk session was not used.

| 😀 Happy | 😐 Neutral | ☹️ Sad | N/A |
|:---:|:---:|:---:|:---:|
| ○ | ○ | ○ | ○ |

## Q3. Would you like to play or talk with EMRA again?

| Yes | Maybe | No |
|:---:|:---:|:---:|
| ○ | ○ | ○ |

---

# F3: Educator Post-Pilot Questionnaire

This questionnaire is completed once by the educator at the end of the pilot.

## Educator Background Information

| Item | Response |
|---|---|
| Center | |
| Role | |
| Years working with autistic children | |
| Number of EMRA sessions run or observed | |
| Language used with EMRA | ○ Arabic only<br>○ Mostly Arabic<br>○ English only<br>○ Mostly English<br>○ About equal Arabic and English |
| Previous experience with educational technology or robots | ○ None<br>○ Some<br>○ Extensive |

## Rating Scale

| Score | Meaning |
|:---:|---|
| 1 | Strongly disagree |
| 2 | Disagree |
| 3 | Neutral |
| 4 | Agree |
| 5 | Strongly agree |

## Section A: Usefulness and Educational Value

| Item | 1 | 2 | 3 | 4 | 5 |
|---|:---:|:---:|:---:|:---:|:---:|
| Q1. EMRA is a useful tool for supporting children’s learning. | ○ | ○ | ○ | ○ | ○ |
| Q2. EMRA’s activities matched the skills we teach. | ○ | ○ | ○ | ○ | ○ |
| Q3. The activities were appropriate for the children. | ○ | ○ | ○ | ○ | ○ |
| Q4. EMRA’s instructions were clear for the children. | ○ | ○ | ○ | ○ | ○ |
| Q5. I would use EMRA again in my sessions. | ○ | ○ | ○ | ○ | ○ |

## Section B: Emotion and Physical Interaction

| Item | 1 | 2 | 3 | 4 | 5 |
|---|:---:|:---:|:---:|:---:|:---:|
| Q6. EMRA responded appropriately when a child appeared upset or frustrated. | ○ | ○ | ○ | ○ | ○ |
| Q7. EMRA’s facial expressions were appropriate to the child’s observed emotional state. | ○ | ○ | ○ | ○ | ○ |
| Q8. EMRA’s gestures and head movements added value to the interaction. | ○ | ○ | ○ | ○ | ○ |
| Q9. The children appeared comfortable interacting with EMRA. | ○ | ○ | ○ | ○ | ○ |
| Q10. EMRA’s movements appeared safe and comfortable for the children. | ○ | ○ | ○ | ○ | ○ |

## Section C: Language

| Item | 1 | 2 | 3 | 4 | 5 |
|---|:---:|:---:|:---:|:---:|:---:|
| Q11. EMRA understood the children’s speech well enough for interaction. | ○ | ○ | ○ | ○ | ○ |
| Q12. EMRA’s voice was clear and easy to understand. | ○ | ○ | ○ | ○ | ○ |
| Q13. EMRA’s Arabic was natural and appropriate for children in this region. | ○ | ○ | ○ | ○ | ○ |
| Q14. EMRA used the correct language for each child. | ○ | ○ | ○ | ○ | ○ |

## Section D: Practical Feasibility

| Item | 1 | 2 | 3 | 4 | 5 |
|---|:---:|:---:|:---:|:---:|:---:|
| Q15. Preparing and starting EMRA was easy. | ○ | ○ | ○ | ○ | ○ |
| Q16. An EMRA session fitted within the available session time. | ○ | ○ | ○ | ○ | ○ |
| Q17. I could use EMRA without needing substantial additional training. | ○ | ○ | ○ | ○ | ○ |
| Q18. EMRA could be used regularly in this center. | ○ | ○ | ○ | ○ | ○ |
| Q19. I could stop or interrupt EMRA easily when needed. | ○ | ○ | ○ | ○ | ○ |
| Q20. I would recommend EMRA for use in another autism center. | ○ | ○ | ○ | ○ | ○ |

## Section E: Open Questions

### 1. What worked best about EMRA?

<br><br><br><br>

### 2. What was the main problem or difficulty when using EMRA?

<br><br><br><br>

### 3. What would you change or improve before EMRA is used regularly?

<br><br><br><br>

---

# Automatically Logged Fields

EMRA automatically records educational, conversational, emotion, action, and system-performance data throughout the pilot.

| Category | Logged Fields |
|---|---|
| **Educational activities** | Item attempted, correctness, response time, attempts, level completed, progress, and help requests |
| **Conversation and emotion** | Transcript, text emotion, vision emotion, fused emotion, channel agreement, reply text, and language routing |
| **Embodied and visual interaction** | Action selected, action source, VLM activation, visual description, and interrupt use |
| **System performance** | End-to-end latency, component latencies, empty transcripts, failures, restarts, and dropped turns |

## Derived Measures for Reporting

| Measure | Derivation |
|---|---|
| **Accuracy per activity** | Correct answers divided by items attempted |
| **Independence rate** | Proportion of activities completed at prompt level 0 |
| **Mean prompt level** | Average of the highest prompt level recorded per activity within a session |
| **Initiation frequency** | Frequency rating of child initiations recorded in F1 Part 2 |
| **Planned-activity completion rate** | Proportion of planned activities recorded as Completed = Yes in F1 Part 1 |
| **Emotion channel agreement rate** | Turns with agreeing text and vision channels divided by turns with a valid visual prediction |
| **Safety override rate** | Turns in which the safety layer replaced the model-selected action divided by all turns containing an action |
| **Recognition failure rate** | Empty transcripts divided by all recorded utterances |
| **Visual grounding rate** | Turns in which the vision-language model was activated divided by all conversational turns |

---

# Privacy

The evaluation uses anonymous child identifiers. Child names and other personally identifying information should not be entered into these forms or published in this repository.

# Related Repository

The EMRA system implementation, educational activities, and interfaces are available in the main repository:

[EMRA Educational Robot](https://github.com/fatimah48/EMRA_Educational-Robot-)

# Research Context

These instruments were developed for the research project:

> **Design and Evaluation of a Multimodal, Emotion-Aware and Embodied Educational Robot for Children with Autism Spectrum Disorder**

King Fahd University of Petroleum and Minerals (KFUPM).

# Citation

If you use or adapt these evaluation instruments, please cite the associated EMRA research publication or thesis. Complete citation details will be added after publication.
